\section{An explicit budget for exponent improvement}
\label{sec:area_law_bootstrap_round_bound}

The finite iteration in the proof of Proposition~9.5
\cite[lines~788--803]{OpenAI2026AreaLaw} contracts an initial exponent
$e_0$ by a factor $1-\ell$ at each round. The following auxiliary estimates
bound the number of rounds in terms of the cyclic-color initial recurrence
parameters of Section~\ref{sec:area_law_initial_exponent_gain}.

\begin{theorem}[Logarithmic round bound]
    \label{thm:area_law_bootstrap_round_log_bound}
    \lean{TNLean.PEPS.AreaLaw.Scan.boxBootstrapRounds_le_ceil_log_div}
    \leanok
    Let $0<e_0\le1$, $0<\varepsilon\le1$ and $0<\ell<1$. Then
    \begin{align}
        m:=\left\lceil\frac{\log(\varepsilon/e_0)}{\log(1-\ell)}\right\rceil_+
        \le\left\lceil\frac{\log(1/\varepsilon)}{\ell}\right\rceil_+.
    \end{align}
\end{theorem}
\begin{proof}\leanok
    Since $\log(1-\ell)\le-\ell<0$ and $\log(1/\varepsilon)\ge0$,
    \begin{align}
        \frac{\log(1/\varepsilon)}{\ell}\log(1-\ell)
        \le-\log(1/\varepsilon)=\log\varepsilon
        \le\log(\varepsilon/e_0).
    \end{align}
    Divide by the negative logarithm and take nonnegative integer ceilings.
\end{proof}

\begin{definition}[Initial round budget]
    \label{def:area_law_initial_round_budget}
    \lean{TNLean.PEPS.AreaLaw.Scan.initialBoxBootstrapRoundBound}
    \leanok
    \uses{thm:area_law_initial_exponent_gain_bounds}
    For the cyclic-color recurrence probability $p$ and scale $M$, set
    \begin{align}
        B:=\left\lceil\frac{400000\log M\log500000}{p}\right\rceil_+.
    \end{align}
    These are the recurrence constants specified above, rather than the
    manuscript's literal numerical choices.
\end{definition}

\begin{theorem}[Bound on the physical bootstrap count]
    \label{thm:area_law_initial_round_budget_bound}
    \lean{TNLean.PEPS.AreaLaw.Scan.boxBootstrapRounds_initial_le}
    \leanok
    \uses{def:area_law_initial_round_budget}
    For the initial exponent $0<e_0<1$, let
    $\ell=(1-e_0)/200000$ and $\varepsilon=2\cdot10^{-6}$.
    Then the round count $m$ above satisfies $m\le B$.
\end{theorem}
\begin{proof}\leanok
    \uses{thm:area_law_bootstrap_round_log_bound,
        thm:area_law_initial_exponent_gain_bounds}
    The gain bound gives
    \begin{align}
        \ell\ge\frac{p}{400000\log M}>0,
        \qquad
        \frac{\log(1/\varepsilon)}{\ell}
        \le\frac{400000\log M\log500000}{p}.
    \end{align}
    Apply the logarithmic round bound and monotonicity of the ceiling.
\end{proof}

\begin{theorem}[The round budget reaches the safe-box exponent]
    \label{thm:area_law_initial_round_budget_reaches}
    \lean{TNLean.PEPS.AreaLaw.Scan.pow_initialBoxBootstrapRoundBound_mul_le}
    \leanok
    \uses{def:area_law_initial_round_budget}
    With the same initial parameters,
    \begin{align}
        (1-\ell)^B e_0\le2\cdot10^{-6}.
    \end{align}
\end{theorem}
\begin{proof}\leanok
    \lean{TNLean.PEPS.AreaLaw.Scan.pow_boxBootstrapRounds_mul_le}
    \uses{thm:area_law_initial_round_budget_bound}
    Since $0<1-\ell<1$ and $m\le B$,
    $(1-\ell)^B e_0\le(1-\ell)^m e_0\le2\cdot10^{-6}$.
\end{proof}
